\scanpage{35}
\mspage{17}
sous $\mathbb F_2^J$, que je vais désigner par $P_{M,\Delta}$

\[
 F(\Delta_M)=\underbrace{P_{M,\Delta}}_{\text{ens des faces de }\Delta}
 \quad\text{torseur sous }\mathbb F_2^J\qquad(J=M_0^*) \tag*{(34)}
\]

Or on a

\[
 \mathbb F_2^J\simeq\underset{\substack{\rotatebox[origin=c]{90}{$\scriptstyle\simeq$}\\ M_0}}{V_J(\mathbb F_2)}\times
 \underset{\text{diag}}{\underline{\mathbb F_2}}\tag*{(35)}
\]

et l'opération sur $P_{M,\Delta}$ de l'élément non nul \deleted{\unsure{$\mathbb F_2$}} du facteur $\mathbb F_2$ n'est autre que l'antipodisme. Ceci posé, on a un iso. canonique

\[
 P_{M,\Delta}=F(\Delta_M)\xleftarrow{\sim}P_M/\mathfrak z
 \qquad\underset{\substack{\|\\\text{Centre de }\mathrm{Gl}(M)}}{\mathfrak z}
 =\{\pm\mathrm{id}_M\}\tag*{(36)}
\]

provenant de l’homomorphisme évident

\diagram{\scriptsize\begin{tikzcd}[ampersand replacement=\&,column sep=0.4em]
{\unsure{(|c}}\&P_M\arrow[dd]\&{= \{(\xi_i)_{i\in J}\in M^{*J}\mid \xi_i^0=i\deleted{\unsure{,|k|*}}\ \forall i\in J,\ \sum\xi_i=0\}}\\
{(37)}\&{}\& {(\xi_i)\longmapsto\{(\dot\xi_i)\}\quad\text{où }\dot\xi=\Lambda^*\xi=\pm\xi\in\Delta\text{ pour }\xi\in M^*}\\
\&F(\Delta_M)\&{= \{f\in\mathcal P_3(\Delta_M)\mid \underbrace{f\to J}_{\substack{\text{induit par}\\\Delta\to J}}\text{ bijectif}\}.}
\end{tikzcd}}

Ceci nous montre à nouveau que $\mathrm{Aut}(\Delta_M)\simeq\mathrm{Gl}(M)'$, de façon précise on a l'iso. de suites exactes

\diagram{\scriptsize\begin{tikzcd}[ampersand replacement=\&,column sep=0.35em]
{(38)}\&1\arrow[r]\&\overbrace{\mathrm{End}(M_0)/\mathbb F_2\cdot\mathrm{id}_{M_0}}^{M_0\times\mathbb F_2}\arrow[r]\arrow[d,"\sim"',"{\ast}"]\&\mathrm{Gl}(M)'\arrow[r]\arrow[d,"\sim"']\&\overbrace{\mathrm{Gl}_{\mathbb F_2}(M_0)}^{\mathfrak S_J}\arrow[r]\arrow[d,"\shortparallel"]\&1\\
\&1\arrow[r]\&\underbrace{\mathbb F_2^J}_{\simeq M_0\times\mathbb F_2}\arrow[r]\&\mathrm{Aut}_{\mathrm{oct}}(\Delta_M)\arrow[r]\&\mathfrak S_J\arrow[r]\&1
\end{tikzcd}}

% slanted
\reversemarginpar\margin{\raggedright\uline{NB} La donnée d'un octaèdre (pas orienté) $\Delta$ équivaut à celle d'un couple $(I,\underline\mu)$, où $I$ est un ens. à quatre éléments (l'ens de ses bifaces -- qui définit l'oct. \uline{orienté} correspondant), et où $\mu=\mathrm{or}(\Delta)$ est un ens à deux éléments ; -- mais on \unsure{constate|voit} que (\unsure{puisq.|puisque} $I$ \unsure{donné|de Kant}) la donnée de $\mu$ équivaut à celle \unsure{du|des} \unsure{c/h|cube} $\unsure{C_\Delta|E_\Delta|I_\Delta}$ (cf plus bas), qui est $\simeq\mu\wedge\underline\omega_I$}\normalmarginpar

A un octaèdre $(\Delta,\underline a_\Delta)$ ou $(\Delta\to J)$ sont associés trois ens. à deux éléments remarquables, correspondant aux trois caractères \inserted{non-triviaux} de $\mathrm{Aut}_{\mathrm{oct}}(\Delta)$, i.e. aux trois hom \deleted{\unsure{non}} nuls de \deleted{$\mathrm{Aut}_{\mathrm{oct}}(\Delta)$}

\[
 \mathrm{Aut}_{\mathrm{oct}}(\Delta)_{\mathrm{ab}}\simeq
 (\mathrm{Aut}_{\mathrm{oct}}^+(\Delta)\times\{1,\underline a_\Delta\})_{\mathrm{ab}}
 \simeq\mathbb F_2\times\mathbb F_2
\]

Ce sont respectivement les ensembles
% archive sheet number: 34
